\documentclass[12pt,a4paper,oneside]{report}
\usepackage[margin=25mm,headheight=20pt]{geometry}
\usepackage{fontspec}
\setmainfont{THSarabunPSK}[Path=assets/,Extension=.ttf,UprightFont=*-Regular,BoldFont=*-Bold,ItalicFont=*-Italic,BoldItalicFont=*-BoldItalic]
\setsansfont{THSarabunPSK}[Path=assets/,Extension=.ttf,UprightFont=*-Regular,BoldFont=*-Bold]
\setmonofont{lmmono10-regular.otf}[Scale=MatchLowercase]
% TH SarabunPSK: body 16pt, titles 18pt, small (header, diagram labels) 14pt
% sizes are Word/PDF points (bp): 16bp = 16.06 TeX pt
\makeatletter
\renewcommand\normalsize{\@setfontsize\normalsize{16.06}{20}%
  \abovedisplayskip 10\p@ \@plus3\p@ \@minus6\p@ \belowdisplayskip \abovedisplayskip
  \abovedisplayshortskip \z@ \@plus3\p@ \belowdisplayshortskip 6\p@ \@plus3\p@ \@minus3\p@}
\renewcommand\small{\@setfontsize\small{14.05}{17}}
\makeatother
\newcommand{\titlesize}{\fontsize{18.07}{22}\selectfont}
\normalsize
\XeTeXlinebreaklocale "th"
\XeTeXlinebreakskip=0pt plus 1pt
\usepackage{pifont,capt-of,graphicx,booktabs,tabularx,longtable,array,amsmath,enumitem,xcolor,tikz,fancyhdr,titlesec,float}
\usetikzlibrary{arrows.meta,positioning,shapes.geometric}
\usepackage[unicode,colorlinks=true,linkcolor=black,urlcolor=black,citecolor=black]{hyperref}
\usepackage{bookmark}
\definecolor{navy}{HTML}{000000}
\definecolor{pale}{HTML}{F2F2F2}
\renewcommand{\contentsname}{สารบัญ}
\renewcommand{\listtablename}{สารบัญตาราง}
\renewcommand{\listfigurename}{สารบัญภาพ}
\renewcommand{\tablename}{ตารางที่}
\renewcommand{\figurename}{ภาพที่}
\renewcommand{\bibname}{บรรณานุกรม}
\renewcommand{\chaptername}{บทที่}
\renewcommand{\appendixname}{ภาคผนวก}
\titleformat{\chapter}[display]{\bfseries\color{navy}\titlesize}{\chaptername\ \thechapter}{5pt}{\titlesize}
\titlespacing*{\chapter}{0pt}{0pt}{22pt}
\titleformat{\section}{\normalsize\bfseries\color{navy}}{\thesection}{0.6em}{}
\titleformat{\subsection}{\normalsize\bfseries}{\thesubsection}{0.6em}{}
\setlength{\parindent}{1.3em}
\setlength{\parskip}{5pt}
\linespread{1}
\setlength{\emergencystretch}{3em}
\tolerance=3000
\setlist{itemsep=3pt,topsep=5pt,leftmargin=1.6em}
\renewcommand{\arraystretch}{1.35}
\newcolumntype{Y}{>{\raggedright\arraybackslash}X}
\pagestyle{fancy}
\fancyhf{}
\fancyhead[L]{\small HENJAI | เห็นใจ}
\fancyhead[R]{\small ทีม Rushhour}
\fancyfoot[C]{\normalsize\thepage}
\renewcommand{\headrulewidth}{0.3pt}
\fancypagestyle{plain}{\fancyhf{}\fancyfoot[C]{\normalsize\thepage}\renewcommand{\headrulewidth}{0pt}}
\DeclareRobustCommand{\file}[1]{\nolinkurl{#1}}
\newcommand{\note}[1]{\par\noindent\colorbox{pale}{\parbox{\dimexpr\linewidth-2\fboxsep}{#1}}\par}
\newcommand{\frontchapter}[1]{\chapter*{#1}\addcontentsline{toc}{chapter}{#1}}
\newcommand{\projectcover}[1]{%
\hypersetup{pageanchor=false}\begin{titlepage}\centering\titlesize\bfseries
HENJAI\par
\vspace{4mm}เห็นใจ\par
\vspace{9mm}#1\par
\vspace{7mm}โครงการผู้ช่วยบรรยายภาพด้วยปัญญาประดิษฐ์\\ที่เข้าใจอารมณ์และบริบทสำหรับผู้พิการทางสายตา\par
\vspace{3mm}Emotion \& Context-Aware AI Vision Assistant\par
\vfill
ผู้จัดทำ\par\vspace{4mm}
นายกิตติพงศ์ สมสอาด\\
นายชินพัฒน์ ขุนเจริญ\\
นายเรข์ณพัศ สุวรรณปริสุทธิ์\par
\vspace{4mm}ทีม Rushhour\\นิสิตชั้นปีที่ 2 หลักสูตรวิศวกรรมคอมพิวเตอร์และเทคโนโลยีดิจิทัล (CEDT)\\ภาควิชาวิศวกรรมคอมพิวเตอร์\\คณะวิศวกรรมศาสตร์ จุฬาลงกรณ์มหาวิทยาลัย\par
\end{titlepage}\hypersetup{pageanchor=true}}

\hypersetup{pdftitle={HENJAI รายงานฉบับสมบูรณ์},pdfauthor={ทีม Rushhour}}
\begin{document}
\projectcover{รายงานฉบับสมบูรณ์}
\pagenumbering{roman}
\frontchapter{บทสรุปผู้บริหาร}
HENJAI หรือ ``เห็นใจ'' เป็นโครงการพัฒนาต้นแบบผู้ช่วยบรรยายภาพภาษาไทยสำหรับผู้พิการทางสายตา โดยมุ่งให้ผู้ใช้เข้าถึงข้อมูลเกี่ยวกับบุคคล สีหน้าที่สังเกตได้ ท่าทาง สภาพแวดล้อม และสิ่งกีดขวางผ่านเสียงพูด แนวคิดหลักคือการลดภาระจากการตีความภาพด้วยตนเอง และให้ข้อมูลที่สัมพันธ์กับตำแหน่งของผู้ใช้ เช่น ``มีเก้าอี้ขวางทางซ้ายของคุณ'' แทนการบรรยายองค์ประกอบภาพอย่างกว้าง ๆ

ระบบที่ส่งมอบประกอบด้วยสองรูปแบบ ได้แก่ ต้นแบบบนเครื่องคอมพิวเตอร์ใน \file{poc/} และต้นแบบเว็บใน \file{web/} ทั้งสองรูปแบบส่งภาพและบริบทระยะสั้นไปยัง Gemini เพื่อสร้างข้อความบรรยาย ส่วนต้นแบบบนเครื่องมีเส้นทางประมวลผลเพิ่มเติมสำหรับตรวจจับบุคคล ท่าทาง ใบหน้า วัตถุ และประมาณความลึก โดยแยกการเตือนระยะใกล้ออกจากการบรรยายความหมายบน\mbox{คลาวด์} ระบบเว็บใช้กล้องและเสียงสังเคราะห์ของเบราว์เซอร์ และไม่มีเส้นทางโมเดลในเครื่องเช่นเดียวกับต้นแบบ Python

การตรวจหลักฐานการทดลองเดิมพบผลลัพธ์ 24 รายการ แบ่งเป็นกลุ่มอารมณ์ ท่าทาง สภาพแวดล้อม และฉากที่ไม่เน้นบุคคล กลุ่มละ 6 ภาพ ในผลชุดปัจจุบันประมวลผลสำเร็จ 21 รายการ คิดเป็นร้อยละ 87.5 และมี 3 รายการผิดพลาดจากโควตา API เวลาเฉลี่ยเฉพาะรายการสำเร็จเท่ากับ 2.243 วินาทีต่อภาพ ค่ามัธยฐาน 2.1 วินาที ชุดก่อนปรับปรุงสำเร็จ 24 รายการ และมีเวลาเฉลี่ย 2.842 วินาที ตัวเลขนี้เป็นการสรุปไฟล์ผลเดิม ไม่ใช่การทดลองควบคุมซ้ำในวันจัดทำรายงาน

เมื่อเทียบกลุ่มอารมณ์ 6 ภาพกับป้ายกำกับชั่วคราวจากคำค้นและปรับชื่อคลาสให้สอดคล้องกัน ผลชุดปัจจุบันตรงกัน 5 จาก 6 ภาพ (หากเทียบข้อความคลาสตรงตัวตามสคริปต์เดิมจะได้ 4 จาก 6) อย่างไรก็ตาม คำค้นไม่ใช่คำตอบมาตรฐานที่ผู้ประเมินตรวจภาพจริง และตารางประเมินคุณภาพโดยมนุษย์ยังไม่มีคะแนน จึงยังไม่สามารถสรุปความแม่นยำทั่วไป อัตราการสร้างข้อมูลที่ไม่มีในภาพ ความพึงพอใจ หรือความปลอดภัยในการเดินทางจริงได้

ผลที่ได้แสดงความเป็นไปได้ของการเชื่อมกล้อง การบรรยายภาพ ความจำระยะสั้น และเสียงภาษาไทยเข้าด้วยกัน ขั้นตอนต่อไปคือจัดทำข้อมูลอ้างอิงที่ตรวจสอบโดยมนุษย์ ประเมินกับผู้ใช้เป้าหมายภายใต้การดูแล และวัดเวลาตั้งแต่รับภาพจนผู้ใช้ได้ยินเสียงจริง รายงานนี้ส่งมอบพร้อมคู่มือผู้ใช้ คู่มือพัฒนา ซอร์สต้นฉบับ ชุดซอร์สแบบ ZIP และข้อมูลทดสอบที่ตรวจย้อนกลับได้

\frontchapter{กิตติกรรมประกาศ}
คณะผู้จัดทำขอขอบคุณคณะวิศวกรรมศาสตร์ จุฬาลงกรณ์มหาวิทยาลัย และหลักสูตรวิศวกรรมคอมพิวเตอร์และเทคโนโลยีดิจิทัล ภาควิชาวิศวกรรมคอมพิวเตอร์ ซึ่งเป็นสังกัดของคณะผู้จัดทำ และขอขอบคุณผู้จัดโครงการฝึกอบรมเทคโนโลยีและปัญญาประดิษฐ์เพื่อพัฒนานวัตกรรมสำหรับผู้พิการ 5 หมวดหมู่ ได้แก่ มหาวิทยาลัยมหิดล ร่วมกับ กสทช. และ สวทช. ที่เปิดโอกาสให้คณะผู้จัดทำนำเสนอแนวคิดและต้นแบบ

ตามข้อมูลกิจกรรมที่คณะผู้จัดทำให้ไว้ ทีมได้รับรางวัลรองชนะเลิศอันดับ 1 จากกิจกรรมเมื่อวันที่ 12--13 กันยายน 2569 ณ โรงแรมรอยัล ริเวอร์ กรุงเทพมหานคร ข้อมูลดังกล่าวเป็นประวัติการเข้าร่วมกิจกรรมของโครงการ และแยกจากผลการประเมินประสิทธิภาพทางเทคนิคที่รายงานในบทที่ 4

ขอขอบคุณผู้เผยแพร่ซอฟต์แวร์ เครื่องมือ แบบจำลอง และภาพทดสอบที่มีแหล่งที่มาในบัญชีข้อมูลของโครงการ องค์ประกอบเหล่านี้ช่วยให้สามารถพัฒนาและตรวจสอบแนวคิดได้อย่างเป็นระบบ การนำทรัพยากรไปใช้ต่อควรคงการให้เครดิตและตรวจเงื่อนไขสิทธิ์ของแต่ละรายการ

\vspace{12mm}\begin{flushright}คณะผู้จัดทำ ทีม Rushhour\\20 ตุลาคม 2569\end{flushright}
\clearpage\tableofcontents
\clearpage\phantomsection\addcontentsline{toc}{chapter}{สารบัญตาราง}\listoftables
\clearpage\phantomsection\addcontentsline{toc}{chapter}{สารบัญภาพ}\listoffigures
\clearpage\pagenumbering{arabic}
\chapter{บทนำ}
\section{ความเป็นมาและความสำคัญของปัญหา}
การรับรู้สิ่งแวดล้อมผ่านการมองเห็นครอบคลุมทั้งตำแหน่งวัตถุ การเปลี่ยนแปลงของสถานการณ์ และสัญญาณทางสังคม ผู้พิการทางสายตาอาจเข้าถึงข้อมูลเหล่านี้ได้ไม่ครบถ้วนจากเสียงแวดล้อมเพียงอย่างเดียว ตัวอย่างเช่น การทราบว่ามีคนอยู่ตรงหน้ายังไม่เพียงพอต่อการเข้าใจว่าอีกฝ่ายกำลังหันมาพูดคุย โบกมือ หรือแสดงสีหน้าที่ควรระมัดระวัง ขณะเดียวกัน การบรรยายภาพทุกส่วนอย่างยืดยาวอาจรบกวนการฟังเสียงรอบตัวและทำให้ข้อมูลสำคัญมาถึงช้า

โครงการจึงเสนอผู้ช่วยที่ใช้กล้องเป็นช่องทางรับข้อมูลและแปลงผลวิเคราะห์เป็นเสียงภาษาไทย โดยให้ความสำคัญกับสิ่งที่เกี่ยวข้องกับผู้ใช้ ณ ขณะนั้น การออกแบบต้องพิจารณาทั้งคุณภาพคำบรรยาย ความกระชับ ความหน่วง การพูดซ้ำ และภาระการควบคุมหน้าจอ ไม่ใช่เพียงจำนวนวัตถุที่โมเดลระบุได้

เอกสารแนวคิดและต้นแบบเดิมใช้ชื่อ Emotion \& Context-Aware AI Vision Assistant และพัฒนาเป็น HENJAI แนวทางที่ปรากฏในโค้ดมีการเปลี่ยนจากต้นแบบบรรยายภาพแบบเดี่ยวไปสู่ระบบที่แบ่งงานระหว่างการตรวจเหตุการณ์ในเครื่องกับการสร้างความหมายบนคลาวด์ เพื่อให้การตอบสนองเหมาะกับลักษณะงานมากขึ้น

\section{วัตถุประสงค์}
\begin{enumerate}
\item พัฒนาต้นแบบรับภาพจากกล้อง วิเคราะห์บุคคล สีหน้า ท่าทาง และบริบท แล้วส่งผลเป็นคำบรรยายและเสียงภาษาไทย
\item ออกแบบคำบรรยายที่สั้นและสัมพันธ์กับตำแหน่งผู้ใช้ โดยให้ความสำคัญกับสิ่งที่ต้องระวังและสิ่งที่เปลี่ยนไป
\item ลดการเรียกบริการเมื่อภาพนิ่ง และลดการพูดซ้ำด้วยประวัติคำบรรยายระยะสั้น
\item ศึกษาความเป็นไปได้ของการแยกงานประมวลผลในเครื่องและบนคลาวด์ พร้อมสรุปผลทดสอบและข้อจำกัดจากหลักฐานที่มี
\item จัดทำเอกสารและไฟล์สำหรับตรวจสอบ ใช้งาน และพัฒนาต่ออย่างเป็นระบบ
\end{enumerate}

\section{ขอบเขตการศึกษา}
ขอบเขตซอฟต์แวร์ประกอบด้วยต้นแบบ Python/Flask สำหรับเครื่องคอมพิวเตอร์ และต้นแบบเว็บที่ใช้ JavaScript กับฟังก์ชันฝั่งเซิร์ฟเวอร์ การวิเคราะห์ความหมายของภาพใช้โมเดลที่มีอยู่ผ่าน API โดยไม่ได้ฝึกโมเดลภาษาภาพหลักขึ้นใหม่ ส่วนการประมวลผลในเครื่องอาศัยโมเดลสำเร็จรูปและกฎที่เขียนในโครงการ

ขอบเขตข้อมูลประเมินอ้างอิง \file{poc/eval/results.json}, \file{results_before.json}, \file{score.csv} และ \file{manifest.json} ภาพที่มีอยู่จริงมี 25 ภาพ แต่ผลการทดลองที่บันทึกไว้ครอบคลุม 24 ภาพแรกเท่านั้น ภาพ \file{emotion_6.jpg} จึงนับเป็นข้อมูลเพิ่มเติมที่ยังไม่มีผลในชุดรายงานนี้ ไม่ขยายตัวหารของผลเดิมเป็น 25

การศึกษานี้ยังไม่ครอบคลุมการรับรองการใช้งานด้านความปลอดภัย การนำทางอิสระ การทดสอบภาคสนามกับผู้พิการทางสายตาอย่างเป็นระบบ หรือการรับรองการเข้าถึงตามมาตรฐานใด การอ่านเอกสารเป็นความสามารถที่ปรากฏในคำสั่งและโค้ดปัจจุบัน แต่ไม่มีชุดทดสอบ OCR ที่เป็นหลักฐานผลสำเร็จ

\section{ประโยชน์ที่คาดว่าจะได้รับ}
ผู้ใช้เป้าหมายอาจเข้าถึงข้อมูลของบุคคลและสภาพแวดล้อมผ่านเสียงได้สะดวกขึ้น ผู้ดูแลและผู้สาธิตมีต้นแบบสำหรับทดลองรูปแบบปฏิสัมพันธ์ ส่วนผู้พัฒนาได้ฐานซอฟต์แวร์และข้อมูลที่สามารถตรวจซ้ำได้ ประโยชน์ต่อผู้ใช้จริงยังเป็นสมมติฐานของการออกแบบที่ต้องประเมินต่อ ไม่ใช่ผลที่ยืนยันจากการทดลองในชุดนี้

\section{นิยามและการอ่านรายงาน}
\begin{table}[H]\centering\caption{คำสำคัญที่ใช้ในรายงาน}
\begin{tabularx}{\linewidth}{lY}\toprule คำ & ความหมายในโครงการ\\\midrule
คำบรรยายบริบท & ข้อความที่เชื่อมสิ่งที่เห็นกับสถานการณ์และตำแหน่งของผู้ใช้\\
สีหน้าหรืออารมณ์ & การตีความสิ่งที่สังเกตได้จากภาพ ไม่ใช่การยืนยันสภาวะภายในบุคคล\\
สายเร็ว & เส้นทางโมเดลในเครื่องของต้นแบบ Python สำหรับข้อมูลใกล้ตัว\\
สายช้า & เส้นทาง Gemini ที่สร้างคำบรรยายและปรับปรุงความจำ\\
ความจำระยะสั้น & คำบรรยายล่าสุดและสรุปบริบทภายใน session\\
หลักฐานผลเดิม & ไฟล์ผลที่บันทึกไว้ก่อนจัดทำรายงาน ไม่ใช่ผลรันใหม่\\\bottomrule
\end{tabularx}\end{table}

\chapter{เอกสารและงานวิจัยที่เกี่ยวข้อง}
\section{การบรรยายภาพด้วยโมเดลภาษาภาพ}
โมเดลภาษาภาพรับข้อมูลภาพร่วมกับคำสั่งข้อความ เพื่อสร้างคำอธิบายที่ตอบโจทย์การใช้งาน เอกสาร Gemini API อธิบายการส่งภาพเป็นข้อมูล inline และการถามตอบเกี่ยวกับภาพ \cite{gemini} โครงการประยุกต์ความสามารถนี้โดยส่งภาพ JPEG ร่วมกับคำสั่งภาษาไทยและประวัติบริบท แล้วขอผลลัพธ์รูปแบบ JSON แทนการนำข้อความอิสระไปใช้โดยตรง

ข้อได้เปรียบคือการรวมข้อมูลบุคคล สถานที่ และข้อความให้อยู่ในคำบรรยายเดียวโดยไม่ต้องฝึกโมเดลแยกทุกงาน ข้อแลกเปลี่ยนของการออกแบบในโครงการคือการพึ่งพาเครือข่าย โควตา และรูปแบบคำตอบจากบริการภายนอก ข้อผิดพลาดดังกล่าวปรากฏจริงในไฟล์ผลทดสอบ จึงต้องแยกความสำเร็จของคำขอออกจากความถูกต้องของเนื้อหา

\section{การตรวจจับวัตถุและการประมาณท่าทาง}
เอกสาร Ultralytics อธิบายงาน pose estimation ว่าให้ตำแหน่งจุดสำคัญของร่างกาย \cite{yolo} โครงการใช้ YOLO11n-Pose สำหรับจุดข้อต่อ และ YOLO11n สำหรับกรอบวัตถุ โดยแปลงผลบางส่วนเป็นกฎท่าทาง เช่น ยกมือหรือกอดอก การแปลงกฎเป็นความหมายทางสังคมเป็นการตัดสินใจของโครงการ ไม่ใช่คุณสมบัติที่ผ่านการรับรองของตัวตรวจจับ

ใน \file{fast.read} โค้ดพิจารณาบุคคลที่มีความสูงของกรอบมากที่สุด และใช้ค่าความเชื่อมั่นของจุดข้อต่อก่อนตีความท่าทาง การเลือกบุคคลด้วยขนาดกรอบเป็นวิธีเชิงประมาณ จึงอาจเลือกคนที่ไม่ใช่คู่สนทนาที่ผู้ใช้ตั้งใจสนใจในฉากหลายคน

\section{การวิเคราะห์ใบหน้าและการจำแนกสีหน้า}
MediaPipe Face Landmarker ให้จุดอ้างอิงใบหน้าและค่าสัญญาณ blendshape \cite{mediapipe} โครงการใช้ผลนี้ประกอบการประมาณทิศทางการมองและสัญญาณใบหน้าในเส้นทางประมวลผลในเครื่อง ส่วนสคริปต์เปรียบเทียบเดิมใช้โมเดล \file{trpakov/vit-face-expression} \cite{fer} ในงานจำแนกสีหน้าหลังตัดบริเวณใบหน้าหรือศีรษะ

ควรแยกการทดลอง FER ดังกล่าวออกจากเส้นทางปัจจุบันของระบบ เพราะ \file{fast.py} ใช้ MediaPipe และความลึก ไม่ได้เรียก ViT classifier ตัวเดียวกับ benchmark เดิม การนำค่าความเร็ว FER จาก README ไปอ้างเป็นเวลาของระบบปัจจุบันทั้งหมดจึงไม่ตรงกับองค์ประกอบที่ทำงานจริง

\section{การประมาณความลึกจากภาพเดียว}
งาน Depth Anything V2 เสนอแบบจำลองประมาณความลึกจากภาพเดียว \cite{depth} โครงการเลือกใช้ checkpoint แบบ metric indoor ผ่าน Transformers และอ่านค่า \file{predicted_depth} เพื่อคำนวณระยะในแถบล่างของภาพ การใช้ค่าทำนายเป็นระยะทางเชิงกายภาพต้องตรวจสอบกับกล้องและสภาพแวดล้อมที่ใช้งานจริง

ในโครงการ พื้นที่แนวตั้งร้อยละ 50--95 ของภาพถูกแบ่งแนวนอนเป็นสามส่วน ใช้ค่ามัธยฐานของแต่ละส่วนแทนระยะทางด้านซ้าย ตรงหน้า และขวา แล้วเลือกส่วนที่ระยะน้อยที่สุด วิธีนี้เป็นกฎสรุปเชิงพื้นที่ ไม่ได้พิสูจน์ว่าตำแหน่งนั้นตรงกับเส้นทางเดินจริงของผู้ใช้ หรือครอบคลุมสิ่งกีดขวางทุกขนาด

\section{กล้องและเสียงบนเว็บ}
เอกสาร API ของเบราว์เซอร์ระบุว่า \file{getUserMedia} ใช้รับสื่อจากอุปกรณ์ภายใต้สิทธิ์ของผู้ใช้และบริบทที่ปลอดภัย \cite{camera} โครงการใช้ API นี้กับกล้องด้านหลังเป็นตัวเลือกที่ต้องการ และย่อภาพก่อนส่งไปยังเซิร์ฟเวอร์ สำหรับเสียง เว็บใช้ \file{speechSynthesis} ซึ่งมีรายการเสียงตามระบบที่ใช้งาน \cite{speech} ดังนั้นการมีเสียงภาษาไทยและพฤติกรรมการเริ่มเล่นเสียงขึ้นกับเบราว์เซอร์และอุปกรณ์

\section{แนวทางการเข้าถึงและตำแหน่งของโครงการ}
WCAG 2.2 เป็นกรอบแนวทางการเข้าถึงเนื้อหาเว็บ \cite{wcag} โครงการนำหลักการลดภาระการควบคุมและให้ข้อมูลผ่านเสียงมาใช้กับหน้าจอปุ่มขนาดใหญ่ แต่ยังไม่มีผลตรวจความสอดคล้องกับ WCAG จึงควรใช้กรอบดังกล่าวเป็นแนวทางประเมินต่อไป

\begin{table}[H]\centering\caption{ความสัมพันธ์ระหว่างองค์ความรู้กับองค์ประกอบของ HENJAI}
\begin{tabularx}{\linewidth}{p{32mm}YY}\toprule องค์ความรู้ & การใช้ในระบบ & ข้อจำกัดที่ต้องทดสอบ\\\midrule
โมเดลภาษาภาพ & บรรยายบริบทเป็นภาษาไทย & ข้อมูลเกินภาพ ความหน่วง โควตา\\
Object/Pose detection & ตำแหน่งวัตถุและข้อต่อ & การบังกัน แสง และมุมกล้อง\\
Face landmarks & สัญญาณใบหน้าและการมอง & หน้าขนาดเล็กหรือหันข้าง\\
Monocular depth & ประมาณความใกล้ในสามทิศ & ความคลาดเคลื่อนระยะจริง\\
Browser media/TTS & รับภาพและสื่อสารด้วยเสียง & สิทธิ์กล้องและเสียงไทย\\\bottomrule
\end{tabularx}\end{table}

ช่องว่างที่โครงการมุ่งแก้คือการรวมข้อมูลหลายชนิดให้เป็นคำบรรยายภาษาไทยที่สัมพันธ์กับผู้ใช้ และลดข้อความซ้ำระหว่างการใช้งานต่อเนื่อง ความแตกต่างนี้เป็นการออกแบบปฏิสัมพันธ์และการเชื่อมระบบ ไม่ใช่ข้ออ้างว่าเสนอโมเดลพื้นฐานใหม่หรือทำได้ดีกว่างานอื่นจากการทดลองเปรียบเทียบมาตรฐาน

\chapter{วิธีดำเนินการ}
\section{ขั้นตอนและแหล่งข้อมูลสำหรับจัดทำรายงาน}
กระบวนการประกอบด้วยการรวบรวมความต้องการจากเอกสารแนวคิด การออกแบบคำบรรยายและหน้าจอ การเชื่อมกล้องกับบริการ AI การเพิ่มการตรวจภาพเปลี่ยนและความจำ จากนั้นจัดทำชุดภาพทดสอบและเก็บผล การจัดทำรายงานฉบับนี้ตรวจเทียบเอกสาร \file{handoff.md}, \file{pitch.md}, \file{slide.md} กับโค้ดปัจจุบันและไฟล์ผล โดยให้โค้ดกับผลที่บันทึกไว้เป็นหลักเมื่อเอกสารเก่าระบุไม่ตรงกัน

ไม่ได้เก็บภาพจากกล้องใหม่ ไม่ได้เรียกโมเดลคลาวด์ซ้ำ และไม่ได้ทดลองกับผู้ใช้เป้าหมายเพิ่มเติมในรอบการจัดทำเอกสาร ดังนั้นคำว่า ``ผลการดำเนินงาน'' ในรายงานครอบคลุมซอฟต์แวร์ที่มีอยู่และการวิเคราะห์หลักฐานเดิม ไม่ใช้แทนผลรับรองระบบพร้อมใช้งานภาคสนาม

\section{การวิเคราะห์ความต้องการ}
\begin{table}[H]\centering\caption{ความต้องการและวิธีรองรับในต้นแบบ}
\begin{tabularx}{\linewidth}{p{13mm}YY}\toprule รหัส & ความต้องการ & วิธีรองรับที่พบในซอร์ส\\\midrule
FR-01 & รับภาพจากสภาพแวดล้อม & OpenCV ใน Python; getUserMedia บนเว็บ\\
FR-02 & บรรยายบุคคลและบริบทเป็นภาษาไทย & Gemini พร้อม SYSTEM\_PROMPT\\
FR-03 & แจ้งข้อมูลที่สัมพันธ์กับผู้ใช้ & กฎคำบรรยายเน้นทิศและสิ่งต้องระวัง\\
FR-04 & ฟังเสียงคำบรรยาย & gTTS/macOS หรือ Browser TTS\\
FR-05 & ลดการพูดเรื่องเดิม & recent 5 ประโยค สรุป 600 ตัวอักษร\\
FR-06 & ลดคำขอเมื่อภาพนิ่ง & คะแนนความแตกต่างของภาพย่อ\\
FR-07 & รองรับการอ่านเอกสาร & reading.is\_document และ reading.text\\
FR-08 & ตรวจข้อมูลระยะใกล้ในเครื่อง & YOLO/Depth/MediaPipe เฉพาะ Python\\\bottomrule
\end{tabularx}\end{table}

ความต้องการด้านคุณภาพประกอบด้วยการเริ่มและหยุดได้ง่าย การหลีกเลี่ยงเสียงทับกัน การแสดงสถานะผิดพลาด การจำกัดบริบทที่ส่งไป และการไม่เปิดเผย API key ต่อเบราว์เซอร์ อย่างไรก็ตาม การมีโค้ดรองรับความต้องการไม่ได้แปลว่าผ่านเกณฑ์คุณภาพแล้ว จำเป็นต้องตรวจประสิทธิผลในสถานการณ์ใช้งานจริงอีกครั้ง

\section{สถาปัตยกรรมระบบ}
\begin{figure}[H]\centering
\begin{tikzpicture}[node distance=9mm,box/.style={draw=navy,fill=pale,rounded corners,align=center,text width=46mm,minimum height=12mm,font=\small},arr/.style={-{Stealth},thick,draw=navy}]
\node[box] (cam) {กล้อง OpenCV / เฟรมล่าสุด};
\node[box,below left=13mm and 1mm of cam] (fast) {ในเครื่อง: YOLO + Depth\\MediaPipe / กฎระยะใกล้};
\node[box,below right=13mm and 1mm of cam] (slow) {ตรวจภาพเปลี่ยน + ความจำ\\Gemini / คำบรรยาย JSON};
\node[box,below=12mm of fast] (say) {ข้อความเตือน + macOS say};
\node[box,below=12mm of slow] (tts) {คำบรรยาย / อ่านเอกสาร\\gTTS หรือ macOS say};
\draw[arr] (cam)--(fast);\draw[arr] (cam)--(slow);\draw[arr] (fast)--(say);\draw[arr] (slow)--(tts);
\end{tikzpicture}
\caption{สถาปัตยกรรมต้นแบบ Python ตามองค์ประกอบในโค้ดปัจจุบัน}
\label{fig:local}
\end{figure}

คลาส \file{Session} ใน \file{poc/core.py} แยกการรับภาพ การวิเคราะห์คลาวด์ และเส้นทางประมวลผลในเครื่องออกจากกัน ข้อมูลเฟรมล่าสุดถูกอ่านภายใต้ lock ส่วน Flask ส่งหน้าจอ สถานะ และวิดีโอ MJPEG ให้เบราว์เซอร์ การตั้งค่าระหว่าง session อยู่ในหน่วยความจำ ยกเว้นค่ากล้องที่เลือกซึ่งมีไฟล์บันทึกเฉพาะเครื่อง

\begin{figure}[H]\centering
\begin{tikzpicture}[node distance=8mm,box/.style={draw=navy,fill=pale,rounded corners,align=center,text width=110mm,minimum height=11mm,font=\small},arr/.style={-{Stealth},thick,draw=navy}]
\node[box] (a) {เบราว์เซอร์: กล้อง / canvas / ตรวจภาพเปลี่ยน};
\node[box,below=of a] (b) {POST /api/analyze: image + recent + summary + mode};
\node[box,below=of b] (c) {Serverless: อ่าน API key จาก environment แล้วเรียก Gemini};
\node[box,below=of c] (d) {JSON กลับสู่เบราว์เซอร์: อัปเดตความจำและ speechSynthesis};
\draw[arr] (a)--(b);\draw[arr] (b)--(c);\draw[arr] (c)--(d);
\end{tikzpicture}
\caption{การไหลของข้อมูลในต้นแบบเว็บ ซึ่งไม่มีชุดโมเดลสายเร็วของ Python}
\end{figure}

ระบบเว็บไม่เก็บความจำ session เป็นฐานข้อมูลฝั่งเซิร์ฟเวอร์ แต่ส่ง \file{recent[]} และ \file{summary} มากับคำขอแต่ละครั้ง คำว่าไม่เก็บ state ในฟังก์ชันนี้ไม่ได้รับประกันว่าผู้ให้บริการภายนอกหรือระบบบันทึกของแพลตฟอร์มไม่มีการเก็บข้อมูล ผู้รับช่วงต้องกำหนดนโยบายการจัดการข้อมูลก่อนเปิดใช้งานจริง

\section{กระบวนการทาง AI บนคลาวด์}
\subsection{การเตรียมภาพและเลือกโหมด}
ต้นแบบ Python แปลงภาพเป็น RGB และย่อความกว้างเมื่อเกิน 768 พิกเซล รักษาสัดส่วน แล้วบันทึกเป็น JPEG คุณภาพ 80 เว็บใช้ canvas กว้าง 768 พิกเซลและ JPEG คุณภาพ 0.8 ภาพสำหรับตรวจการเปลี่ยนแปลงถูกลดเหลือ 64 คูณ 36 พิกเซลในระดับเทา คะแนนความเปลี่ยนแปลงนิยามเป็น
\begin{equation}
D_t=\frac{1}{255N}\sum_{i=1}^{N}|G_t(i)-G_{t-1}(i)|,\qquad N=64\times36.
\end{equation}
ค่าต่ำหมายถึงภาพใกล้เคียงกัน ค่าสูงหมายถึงความต่างมาก แต่การเคลื่อนของกล้องก็ทำให้ค่าสูงได้โดยไม่มีวัตถุใหม่ ดังนั้นคะแนนนี้เป็นตัวควบคุมการเรียก AI ไม่ใช่ตัวตรวจอันตรายโดยตรง

\begin{table}[H]\centering\caption{โหมดคำสั่งและค่าควบคุมการเปลี่ยนภาพ}
\begin{tabularx}{\linewidth}{lYY}\toprule โหมด & การทำงาน & เงื่อนไขหลัก\\\midrule
first & บรรยายภาพรวมครั้งแรก & ยังไม่มีผลสำเร็จก่อนหน้า\\
skip & ข้ามการเรียกบริการ & ภาพเปลี่ยนน้อยกว่า still\\
delta & บอกเฉพาะสิ่งเปลี่ยน & มากกว่า still แต่น้อยกว่า scene\\
scene & บรรยายสถานที่ใหม่ & มากกว่าหรือเท่ากับ scene\\
alert & เริ่มจากสิ่งต้องระวัง & มากกว่าหรือเท่ากับ alert\\\bottomrule
\end{tabularx}\end{table}

เว็บกำหนด still เท่ากับ 0.02, scene เท่ากับ 0.18 และ alert เท่ากับ 0.45 ส่วน Python คำนวณเกณฑ์จาก noise กล้อง ด้วยตัวคูณ 2.5, 9 และ 22 ตามลำดับ ตัวอย่าง noise 0.02 ให้เกณฑ์ 0.05, 0.18 และ 0.44 การสอบเทียบใช้ค่ามัธยฐานความต่างของเฟรมช่วงเริ่มต้น 40 ค่า จึงควรตั้งกล้องให้นิ่งในช่วงเริ่ม session

\subsection{คำสั่งและสัญญาข้อมูล}
คำสั่งหลักระบุให้ตอบ JSON และใช้คลาสอารมณ์จำกัด ได้แก่ happy, sad, angry, fear, surprised, disgust, neutral และ no\_person พร้อมคำอธิบายภาษาไทย กฎสำคัญคือห้ามเปลี่ยนสีหน้าลบเป็นคำบรรยายบวก หากไม่ชัดให้ลดความเชื่อมั่น และให้กล่าวถึงผู้ใช้โดยตรง

\begin{table}[H]\centering\caption{ข้อมูลผลลัพธ์ของการวิเคราะห์ภาพ}
\begin{tabularx}{\linewidth}{p{46mm}Y}\toprule ฟิลด์ & หน้าที่\\\midrule
\file{narrative_th} & คำบรรยายสั้นภาษาไทย โดยคำสั่งตั้งเป้าไม่เกิน 20 คำ\\
\file{emotion} & label, confidence และ detail\_th ของสีหน้าที่ตีความ\\
\file{body_language} & รายละเอียดท่าทางที่สังเกตได้\\
\file{environment} & ฉาก วัตถุ และข้อความเกี่ยวกับสิ่งต้องระวัง\\
\file{reading} & สถานะเอกสารและข้อความที่อ่านได้\\
\file{memory_th} & สรุปบริบทเพื่อใช้รอบถัดไป\\\bottomrule
\end{tabularx}\end{table}

ค่าการเรียกในซอร์สใช้ temperature 0.4 และจำกัดผลลัพธ์ 512 tokens พร้อม timeout 25 วินาที หาก JSON แปลงไม่ได้จะคืนข้อความบางส่วนเป็นคำบรรยาย fallback จึงยังไม่มีการยืนยัน schema อย่างเข้มงวด ข้อกำหนด ``ไม่เกิน 20 คำ'' เป็นข้อกำหนดใน prompt และยังไม่มีตัวนับคำภาษาไทยบังคับหลังรับคำตอบ

\subsection{ความจำ การกันซ้ำ และการอ่านเอกสาร}
Python เก็บข้อความล่าสุด 5 ประโยค และตรวจความคล้ายด้วย \file{difflib.SequenceMatcher} ที่เกณฑ์มากกว่า 0.8 เว็บใช้สัดส่วนคำร่วมที่แยกด้วยช่องว่าง ซึ่งต่างจาก Python และอาจมีข้อจำกัดกับภาษาไทยที่ไม่ได้เว้นวรรคทุกคำ สรุปบริบทจากโมเดลถูกตัดให้ไม่เกิน 600 ตัวอักษรก่อนส่งต่อ

เมื่อโมเดลระบุว่าเป็นเอกสาร ระบบใช้ \file{reading.text} เป็นเนื้อหาสำหรับอ่านต่อจากคำเกริ่น การอ่านเอกสารข้ามการกันข้อความซ้ำแบบคำบรรยายฉาก ข้อจำกัดผลลัพธ์ 512 tokens อาจทำให้เอกสารยาวถูกตัด จึงไม่ควรอ้างว่ารองรับการอ่านเอกสารทุกหน้าได้ครบถ้วนจากโค้ดเพียงอย่างเดียว

\section{กระบวนการประมวลผลในเครื่อง}
เส้นทางปัจจุบันโหลด YOLO11n-Pose, YOLO11n, Depth Anything V2 Metric Indoor Small และ MediaPipe Face Landmarker ใช้ MPS หากมี มิฉะนั้นเลือก CPU โดย MediaPipe ตั้ง delegate เป็น CPU โมเดลความลึกต้องดาวน์โหลดจากผู้ให้บริการเมื่อยังไม่มี cache แม้ขั้นตอน inference จะออกแบบให้ทำในเครื่องก็ตาม

ผลระยะล่างของภาพถูกแปลงเป็น \file{near_m} และทิศ \file{near_side} ก่อนเข้า \file{fast.phrase} กฎเตือนปัจจุบันให้พูดเมื่อระยะต่ำกว่า 0.5 เมตรและสถานะพร้อมเตือน แล้วปิดการเตือนซ้ำจนระยะมากกว่า 0.9 เมตร เกณฑ์ 1.5 เมตรใช้แสดงข้อมูลความใกล้ ไม่ใช่เกณฑ์สั่งพูด อัลกอริทึมยังคำนวณสัญญาณเข้ามาใกล้ แต่ไม่ได้ใช้สัญญาณนั้นเป็นเงื่อนไขเตือนทุกกรณี

\note{เกณฑ์ระยะเป็นค่าทดลองในโค้ด ไม่ใช่ระยะปลอดภัยที่ผ่านการรับรอง และยังไม่มีหลักฐานสอบเทียบระยะจริงครบหลายสภาพแวดล้อมในชุดส่งมอบ}

\section{การออกแบบการทดสอบและการวัดผล}
ภาพ 24 รายการในผลเดิมแบ่งเป็น 4 กลุ่ม กลุ่มละ 6 ภาพ สคริปต์ \file{eval/run.py} วิเคราะห์แต่ละภาพด้วยโหมด first จึงไม่ใช่การทดสอบความจำต่อเนื่อง สคริปต์บันทึกเวลา label ความเชื่อมั่น คำบรรยาย ข้อควรระวัง และข้อผิดพลาด แล้วเตรียม CSV สำหรับผู้ประเมิน

กำหนดความครอบคลุมเป็นจำนวนรายการสำเร็จหารจำนวนรายการทั้งหมด ส่วนเวลาเฉลี่ยนับเฉพาะรายการที่ไม่มีข้อผิดพลาด เกณฑ์สำหรับผู้ประเมินใน CSV ประกอบด้วย correctness แบบ 0/1, helpfulness แบบ 1--5, conciseness แบบ 1--5 และ hallucination แบบ 0/1 ต้องตรวจภาพจริงก่อนกรอก ไม่อนุมานคะแนนจากชื่อไฟล์หรือคำค้น

การเทียบผลก่อนและหลังควรใช้ชุดภาพคู่กันและแยกข้อผิดพลาดบริการออกจากคุณภาพโมเดล ในรายงานนี้เพิ่มสรุปแบบจับคู่ 21 ภาพที่สำเร็จทั้งสองชุด แต่ไม่ตีความว่าเป็นการทดลองควบคุม เนื่องจากไม่มีบันทึกเงื่อนไขบริการและสภาพเครือข่ายครบถ้วน

\chapter{ผลการดำเนินงาน}
\section{ผลลัพธ์ด้านระบบ}
ซอร์สปัจจุบันมีหน้าจอสำหรับผู้ใช้หลัก หน้าจอพัฒนาใน Python การส่งภาพและรับ JSON การใช้ประวัติคำบรรยาย การข้ามภาพนิ่ง และการสร้างเสียงภาษาไทย เส้นทางในเครื่องมีการเตือนระยะชิดตัวตามกฎที่กำหนด ส่วนเว็บใช้การประมวลผลกล้องในเบราว์เซอร์และบริการ Gemini เท่านั้น

\begin{table}[H]\centering\caption{สถานะความสามารถและขอบเขตหลักฐาน}
\begin{tabularx}{\linewidth}{p{39mm}YY}\toprule ความสามารถ & สิ่งที่พบ & หลักฐานที่ยังขาด\\\midrule
กล้องและคำบรรยาย & มีโค้ดทั้งสองแพลตฟอร์ม & รอบทดสอบกล้องปัจจุบัน\\
อารมณ์และบริบท & มี JSON ผลเดิม 21 ภาพสำเร็จ & ป้ายกำกับมนุษย์ที่ยืนยันแล้ว\\
กันพูดซ้ำ/ภาพนิ่ง & มีโค้ดและสคริปต์ทดสอบ & ผลทดสอบการฟังกับผู้ใช้จริง\\
เตือนระยะในเครื่อง & มีโมเดลและกฎ hysteresis & การสอบเทียบและ hazard recall\\
อ่านเอกสาร & มี prompt และเส้นทางอ่านเสียง & ชุด OCR และความครบถ้วนข้อความ\\
การเข้าถึง & มีปุ่มหลักและสถานะเสียง & ผลทดสอบ screen reader/WCAG\\\bottomrule
\end{tabularx}\end{table}

\section{ผลการประมวลผลภาพที่บันทึกไว้}
\begin{table}[H]\centering\caption{สถิติจากผลเดิม แสดงเวลาเฉพาะรายการสำเร็จเป็นวินาที}
\begin{tabularx}{\linewidth}{Yrrrrr}\toprule ชุด & ทั้งหมด & สำเร็จ & ผิดพลาด & เฉลี่ย & มัธยฐาน\\\midrule
ก่อนปรับปรุง & 24 & 24 & 0 & 2.842 & 2.75\\
ชุดปัจจุบัน & 24 & 21 & 3 & 2.243 & 2.10\\\bottomrule\end{tabularx}\end{table}
\begin{table}[H]\centering\caption{ความครอบคลุมและเวลาเฉลี่ยแยกกลุ่มในผลชุดปัจจุบัน}
\begin{tabularx}{\linewidth}{Yrrrr}\toprule กลุ่ม & ทั้งหมด & สำเร็จ & ผิดพลาด & วินาที\\\midrule
อารมณ์ & 6 & 6 & 0 & 2.083\\
ท่าทาง & 6 & 6 & 0 & 2.350\\
สภาพแวดล้อม & 6 & 6 & 0 & 2.333\\
ไม่เน้นบุคคล & 6 & 3 & 3 & 2.167\\\bottomrule\end{tabularx}\end{table}
ข้อผิดพลาดทั้งสามรายการในชุดปัจจุบันคือ \file{nolabel_3.jpg}, \file{nolabel_4.jpg} และ \file{nolabel_5.jpg} ซึ่งบันทึก \file{429 RESOURCE_EXHAUSTED} จึงไม่รวมเวลาในค่าเฉลี่ยของรายการสำเร็จ และไม่แปลงรายการเหล่านี้เป็นคำตอบผิดทางความหมายโดยอัตโนมัติ

การเทียบค่าเฉลี่ยรวมก่อนและหลังมีตัวหารต่างกัน จึงเสี่ยงทำให้เข้าใจว่าความเร็วดีขึ้นเพราะการปรับระบบเพียงอย่างเดียว สำหรับภาพ 21 รายการที่สำเร็จทั้งคู่ ค่าเฉลี่ยก่อนเท่ากับ 2.867 วินาที และหลังเท่ากับ 2.243 วินาที แม้จำนวนภาพเท่ากัน แต่เงื่อนไขบริการภายนอกไม่ได้ถูกตรึงไว้ครบ จึงรายงานเป็นข้อสังเกตจาก log เดิม

\begin{figure}[H]\centering
\begin{tikzpicture}[x=1.8cm,y=0.85cm]
\draw[->] (0,0)--(3.5,0) node[right,font=\small]{วินาที};
\foreach \x in {0,1,2,3}{\draw (\x,0)--(\x,-.08) node[below,font=\small]{\x};}
\fill[navy!55] (0,.5) rectangle (2.842,1.1);
\fill[navy] (0,1.6) rectangle (2.243,2.2);
\node[left,font=\small] at (0,.8){ก่อน: 24 ภาพ};
\node[left,font=\small] at (0,1.9){หลัง: 21 ภาพ};
\node[right,font=\small] at (2.842,.8){2.842};
\node[right,font=\small] at (2.243,1.9){2.243};
\end{tikzpicture}
\caption{เวลาเฉลี่ยเฉพาะคำขอสำเร็จจากผลเดิม ตัวหารสองชุดต่างกัน จึงไม่ใช่ข้อสรุปเชิงเหตุและผล}
\end{figure}

\section{การวิเคราะห์กลุ่มอารมณ์รายภาพ}
\begin{table}[H]\centering\caption{ผลกลุ่มอารมณ์เทียบกับป้ายกำกับชั่วคราวจากคำค้น}
\begin{tabularx}{\linewidth}{lYYY}\toprule ภาพ & ป้ายชั่วคราว & ผลชุดปัจจุบัน & การตรงกัน\\\midrule
emotion\_0 & happy & happy & ตรง\\
emotion\_1 & sad & neutral & ไม่ตรง\\
emotion\_2 & angry & angry & ตรง\\
emotion\_3 & surprise & surprised & ตรงหลังปรับชื่อ\\
emotion\_4 & sad & sad & ตรง\\
emotion\_5 & happy & happy & ตรง\\\bottomrule
\end{tabularx}\end{table}

ตารางข้างต้นนับได้ 5 จาก 6 หลังปรับชื่อ surprise/surprised ให้ตรงกัน ส่วนการเทียบข้อความคลาสตรงตัวกับ \file{TRUTH} ใน \file{bench_local.py} จะนับได้ 4 จาก 6 เพราะคลาส surprised ไม่เท่ากับ surprise ทั้งสองตัวเลขวัดคนละวิธี จึงไม่ควรใช้คำว่า ``ความแม่นยำ 4/6'' โดยไม่ระบุวิธีปรับชื่อ รายงานนี้ใช้ 5/6 เฉพาะการจับคู่คลาสที่ปรับชื่อแล้ว และคง 4/6 ไว้เพื่ออธิบายผลสคริปต์เดิม

กรณี \file{emotion_1.jpg} คำค้นระบุการร้องไห้ แต่คำบรรยายในผลปัจจุบันตีความว่าอาจกำลังหาวหรือจาม จึงเป็นกรณีที่ควรให้ผู้ประเมินตรวจภาพและบริบทจริงก่อนตัดสิน ส่วน \file{emotion_2.jpg} เปลี่ยนจากคำตอบเดิม ``playful'' เป็น ``angry'' หลังข้อกำหนดคำสั่งถูกปรับให้รายงานสีหน้าลบตรงไปตรงมา อย่างไรก็ตาม ไม่มีบันทึกการทดลองควบคุมที่แยกผลของ prompt ออกจากปัจจัยอื่นทั้งหมด

\section{คุณภาพคำบรรยายและผลที่ยังประเมินไม่ได้}
ใน \file{score.csv} มี 24 แถว แต่ช่องคะแนนจากมนุษย์ครบทั้งสี่ด้านยังว่างทั้งหมด จึงไม่มีค่าความถูกต้องเชิงเนื้อหา ความเป็นประโยชน์ ความกระชับ หรือ hallucination rate ที่ผ่านการประเมิน ตัวอย่างคำบรรยายที่อ่านดูสมเหตุสมผลไม่เพียงพอจะยืนยันว่าไม่มีข้อมูลเกินภาพ

เอกสาร README เดิมระบุผล FER หลังตัดหัวด้วย YOLO ประมาณ 72 มิลลิวินาที และ YOLO11-Pose ประมาณ 42 มิลลิวินาที แต่ไม่มีตารางผลดิบของ benchmark ดังกล่าวในชุดไฟล์ที่ตรวจพบ อีกทั้งองค์ประกอบสายเร็วปัจจุบันเพิ่มโมเดลความลึกและ MediaPipe จึงไม่นำเวลารวมประมาณ 110 มิลลิวินาทีในเอกสารเก่ามาเป็นผลวัดระบบปัจจุบัน

\section{ข้อจำกัดและการตรวจย้อนกลับ}
ข้อจำกัดสำคัญคือภาพมีจำนวนน้อย การเลือกภาพอาศัยคำค้น ความหมายของป้ายกำกับยังไม่ผ่านผู้ประเมิน และการทดลองเป็นภาพนิ่ง ไม่ใช่วิดีโอต่อเนื่องจากมุมกล้องของผู้ใช้จริง มีความล้มเหลวจากโควตาและไม่มีเวลาตั้งแต่กดเริ่มจนได้ยินเสียงครบวงจร จึงควรตีความผลเป็นหลักฐานความเป็นไปได้ของต้นแบบ

\file{delivery/data/analysis.json} สรุปจำนวนภาพ รายการผิดพลาด สถิติเวลา และวิธีจับคู่คลาส พร้อม \file{case_index.csv} สำหรับเชื่อมภาพกับผล ชุดข้อมูลคงไฟล์ผลเดิมโดยไม่เขียนทับและมีบัญชี SHA-256 เพื่อยืนยันความตรงกันของไฟล์ส่งมอบ

\chapter{สรุปและข้อเสนอแนะ}
\section{สรุปผลการศึกษา}
โครงการได้ต้นแบบที่เชื่อมการรับภาพ การวิเคราะห์บริบทด้วย AI และการสื่อสารด้วยเสียงภาษาไทย ครอบคลุมทั้งการใช้งานบนเครื่องคอมพิวเตอร์และเว็บ การออกแบบคำบรรยายเน้นผู้ใช้เป็นศูนย์กลาง มีการลดคำขอเมื่อภาพนิ่งและจำบริบทระยะสั้น ส่วนเส้นทางในเครื่องแยกการประเมินระยะใกล้ออกจากการสร้างคำบรรยายบนคลาวด์

หลักฐานผลเดิมยืนยันว่ามีการเรียกวิเคราะห์ภาพสำเร็จ 21 จาก 24 รายการในชุดปัจจุบัน โดยเฉลี่ย 2.243 วินาทีต่อรายการสำเร็จ แต่ยังไม่มีคะแนนคุณภาพจากมนุษย์หรือผลภาคสนามที่เพียงพอ จึงบรรลุการพัฒนาต้นแบบและการรวบรวมหลักฐานขั้นต้น ส่วนเป้าหมายด้านคุณภาพชีวิตและความปลอดภัยต้องศึกษาต่อ

\section{การนำไปใช้ประโยชน์}
สามารถใช้ต้นแบบเป็นสื่อสาธิตการประยุกต์โมเดลภาษาภาพกับเทคโนโลยีช่วยเหลือ ใช้สำรวจรูปแบบคำบรรยายที่ผู้ใช้ต้องการ และใช้เป็นฐานสำหรับพัฒนาชุดประเมินภาษาไทย การสาธิตควรอยู่ในสภาพแวดล้อมควบคุม มีผู้ดูแล และแยกการบรรยายสถานการณ์ออกจากการพึ่งพาระบบเพื่อตัดสินใจเดินทาง

ข้อมูลและคู่มือที่ส่งมอบช่วยให้ผู้รับช่วงตรวจว่าคำบรรยายมาจากไฟล์ใด เปลี่ยนโมเดลหรือ prompt อย่างมีหลักฐาน และวางการทดลองครั้งต่อไปโดยไม่ทำลายผลเดิม ไม่ควรนำผลชนะรางวัลหรือการมีต้นแบบไปใช้แทนหลักฐานความปลอดภัยทางเทคนิค

\section{ข้อเสนอแนะตามลำดับความสำคัญ}
\begin{enumerate}
\item จัดทำคำตอบอ้างอิงจากผู้ประเมินอย่างน้อยสองคนต่อภาพ แยกสีหน้าที่เห็นจริงออกจากการคาดเดาอารมณ์ และตกลงแนวทางจัดการภาพกำกวมก่อนสรุปคะแนน
\item วัดความครอบคลุมและความล้มเหลวควบคู่กับคุณภาพ โดยตรึงชุดภาพ เวอร์ชัน prompt โมเดล และค่าการเรียก พร้อมเก็บผลดิบทุกครั้ง
\item วัดเวลา end-to-end รวมการรับภาพ อัปโหลด วิเคราะห์ สังเคราะห์เสียง และการเริ่มได้ยินเสียง แยกค่ากลางกับค่าช่วงช้า และทดสอบเครือข่ายขาดหาย
\item สอบเทียบระยะกับวัตถุจริงหลายชนิด หลายพื้นผิวและระดับแสง รวมกล้องติดศีรษะ ก่อนกำหนดเงื่อนไขเตือนที่เหมาะสม
\item ทดสอบปุ่มหลักและเสียงกับผู้พิการทางสายตาภายใต้ความยินยอมและการดูแล เก็บข้อมูลว่าคำบรรยายช่วยตัดสินใจจริงหรือสร้างภาระการฟัง
\item ปรับการตรวจ schema และข้อมูลรับเข้า เพิ่มการจำกัดคำขอและควบคุมการเข้าถึงก่อนให้บริการสาธารณะ พร้อมกำหนดนโยบายภาพและเสียง
\end{enumerate}

\section{เกณฑ์สำหรับการประเมินรอบถัดไป}
ก่อนเริ่มควรกำหนดจำนวนกรณี สัดส่วนสถานการณ์ เกณฑ์ความผิดพลาดร้ายแรง และเวลาตอบสนองที่ผู้ใช้ยอมรับร่วมกับผู้เกี่ยวข้อง โดยใช้ผลต้นแบบนี้เป็นฐานตั้งสมมติฐาน ไม่เติมค่าเป้าหมายย้อนหลังให้ดูเหมือนผ่านแล้ว การตัดสินพร้อมใช้งานควรอาศัยผลคุณภาพ ความครอบคลุม และประสบการณ์ผู้ใช้ร่วมกัน

\clearpage\addcontentsline{toc}{chapter}{บรรณานุกรม}
\begin{thebibliography}{99}
\bibitem{gemini} Google. \textit{Gemini API: Image understanding}. \url{https://ai.google.dev/gemini-api/docs/image-understanding}. สืบค้น 20 ตุลาคม 2569.
\bibitem{yolo} Ultralytics. \textit{Pose Estimation}. \url{https://docs.ultralytics.com/tasks/pose/}. สืบค้น 20 ตุลาคม 2569.
\bibitem{mediapipe} Google. \textit{Face landmark detection guide}. \url{https://ai.google.dev/edge/mediapipe/solutions/vision/face_landmarker}. สืบค้น 20 ตุลาคม 2569.
\bibitem{depth} Yang, L., Kang, B., Huang, Z., Zhao, Z., Xu, X., Feng, J., and Zhao, H. (2024). \textit{Depth Anything V2}. arXiv:2406.09414. \url{https://arxiv.org/abs/2406.09414}.
\bibitem{fer} trpakov. \textit{vit-face-expression: Model card}. \url{https://huggingface.co/trpakov/vit-face-expression}. สืบค้น 20 ตุลาคม 2569.
\bibitem{camera} MDN Web Docs. \textit{MediaDevices: getUserMedia() method}. \url{https://developer.mozilla.org/en-US/docs/Web/API/MediaDevices/getUserMedia}. สืบค้น 20 ตุลาคม 2569.
\bibitem{speech} MDN Web Docs. \textit{SpeechSynthesis}. \url{https://developer.mozilla.org/en-US/docs/Web/API/SpeechSynthesis}. สืบค้น 20 ตุลาคม 2569.
\bibitem{wcag} W3C. (2024). \textit{Web Content Accessibility Guidelines (WCAG) 2.2}. \url{https://www.w3.org/TR/WCAG22/}. สืบค้น 20 ตุลาคม 2569.
\bibitem{repo} ทีม Rushhour. \textit{ซอร์สโค้ดและบันทึก HENJAI}. \file{poc/}, \file{web/}, \file{handoff.md}, \file{pitch.md}, \file{research.md}, \file{slide.md}. รุ่นซอร์ส \file{e327cd1}; ตรวจไฟล์ 20 ตุลาคม 2569.
\bibitem{eval} ทีม Rushhour. \textit{หลักฐานการทดลองภาพนิ่ง}. \file{poc/eval/manifest.json}, \file{results_before.json}, \file{results.json}, \file{score.csv}. ไฟล์ผลเดิมของโครงการ; ไม่ได้เรียก API ใหม่ในรอบจัดทำรายงาน.
\bibitem{event} ข้อมูลจากคณะผู้จัดทำ. \textit{รายชื่อ สังกัด และผลเข้าร่วมโครงการฝึกอบรมเทคโนโลยีและปัญญาประดิษฐ์เพื่อพัฒนานวัตกรรมสำหรับผู้พิการ 5 หมวดหมู่}. กิจกรรม 12--13 กันยายน 2569; ให้ข้อมูลสำหรับรายงาน 20 ตุลาคม 2569.
\end{thebibliography}
\appendix
\renewcommand{\chaptername}{ภาคผนวก}
\chapter{บัญชีไฟล์ส่งมอบและข้อมูลทดสอบ}
\section{รายการส่งมอบ}
\begin{longtable}{p{43mm}p{105mm}}
\caption{เอกสารและไฟล์ที่จัดเตรียมเพื่อส่งมอบ}\\\toprule รายการ & ตำแหน่งภายใน delivery\\\midrule\endfirsthead
\toprule รายการ & ตำแหน่งภายใน delivery\\\midrule\endhead
รายงานฉบับสมบูรณ์ & \file{pdf/report.pdf} และ \file{latex/report.tex}\\
คู่มือผู้ใช้ & \file{pdf/user-manual.pdf} และ \file{latex/user-manual.tex}\\
คู่มือพัฒนา & \file{pdf/system-manual.pdf} และ \file{latex/system-manual.tex}\\
ไฟล์ต้นฉบับซอฟต์แวร์ & \file{source_files/poc/} และ \file{source_files/web/} พร้อมไฟล์ประกอบที่จำเป็น\\
ซอร์สบีบอัด & \file{henjai-source.zip}\\
ภาพและข้อมูลทดสอบ & \file{data/img/}, \file{data/manifest.json}, ผลเดิม JSON และ CSV\\
สรุปตรวจสอบข้อมูล & \file{data/analysis.json}, \file{data/case_index.csv}, \file{data/credits.csv}\\
หลักฐานไฟล์ & \file{SOURCE_MANIFEST.csv}, \file{DATA_MANIFEST.csv}, \file{SHA256SUMS.txt}\\
แนวทางคอมไพล์ & \file{README.md} และ \file{latex/build.sh}\\\bottomrule
\end{longtable}

ชุดต้นฉบับไม่รวม environment ส่วนตัว API key ประวัติ Git cache หรือไฟล์ตั้งค่ากล้องเฉพาะเครื่อง ทั้งหมดไม่ใช่ซอร์สซอฟต์แวร์ที่จำเป็นต่อการรับช่วง โมเดลที่มีอยู่ในโฟลเดอร์โครงการถูกเก็บไว้กับชุดซอร์ส ส่วนโมเดลที่โค้ดดาวน์โหลดภายหลังต้องตรวจสิทธิ์และเตรียม cache แยกตามคู่มือพัฒนา

\section{โครงสร้างข้อมูลทดสอบและสิทธิ์}
\file{manifest.json} มี 25 รายการพร้อมชื่อไฟล์ กลุ่ม คำค้น ผู้สร้าง และแหล่งที่มา ภาพจากแหล่งออนไลน์มีรหัสสิทธิ์ by หรือ by-sa ตามที่บันทึกไว้ โดยข้อมูลนี้คงจากต้นฉบับ ไม่ใช่ผลตรวจใบอนุญาตใหม่ทุกภาพ ภาพ \file{emotion_6.jpg} ถูกระบุว่าเป็นภาพที่ผู้ใช้จัดหาและยังต้องตรวจสิทธิ์ก่อนเผยแพร่สาธารณะ

\file{credits.csv} สะท้อนบัญชีผู้สร้างและ URL ของทุกภาพ เพื่อให้ผู้รับช่วงตรวจรายละเอียดเวอร์ชันใบอนุญาตและคงเครดิตก่อนเผยแพร่ซ้ำ ชุดส่งมอบนี้ใช้เพื่อรับช่วงและตรวจงานโครงการ ไม่ได้สรุปว่าทรัพยากรบุคคลที่สามทุกชิ้นอนุญาตใช้เชิงพาณิชย์โดยไม่มีเงื่อนไข

\section{บัญชีรายภาพ}
\textbf{สรุปสั้น:} ชุดทดสอบมีภาพ 25 ภาพ แบ่งเป็น 4 กลุ่ม ทดสอบจริง 24 ภาพ ระบบตอบกลับสำเร็จ 21 ภาพ (ร้อยละ 87.5) ใช้เวลาเฉลี่ย 2.24 วินาทีต่อภาพ อีก 3 ภาพไม่ได้ผลเพราะใช้โควตา API หมด ไม่ใช่เพราะระบบตีความภาพผิด

\textbf{วิธีอ่านตาราง:} แต่ละช่องคือภาพหนึ่งภาพ เช่น แถว 0 คอลัมน์อารมณ์คือไฟล์ \file{emotion_0.jpg}
\begin{itemize}
\item \textbf{ตัวเลข} คือจำนวนวินาทีตั้งแต่ส่งภาพจนได้คำบรรยายกลับมา
\item \textbf{\ding{55} โควตา} คือบริการ Gemini ปฏิเสธคำขอเพราะเรียกใช้เกินโควตา (รหัส HTTP 429)
\item \textbf{ไม่ได้ทดสอบ} คือภาพ \file{emotion_6.jpg} ที่ผู้ใช้ส่งมาเพิ่มหลังการทดสอบรอบนี้ จึงยังไม่มีผล
\end{itemize}
\begin{table}[H]\centering
\caption{เวลาที่ระบบใช้ตอบต่อภาพ (วินาที) แยกตามกลุ่มภาพ}
\begin{tabular}{ccccc}\toprule
ภาพที่ & อารมณ์ & ท่าทาง & สภาพแวดล้อม & ไม่เน้นบุคคล\\
 & (emotion) & (body) & (env) & (nolabel)\\\midrule
0 & 2.5 & 2.1 & 2.2 & 2.2\\
1 & 2.0 & 1.9 & 2.3 & 2.4\\
2 & 2.0 & 2.1 & 3.1 & 1.9\\
3 & 2.1 & 2.4 & 2.6 & \ding{55} โควตา\\
4 & 1.8 & 3.2 & 1.9 & \ding{55} โควตา\\
5 & 2.1 & 2.4 & 1.9 & \ding{55} โควตา\\
6 & ไม่ได้ทดสอบ & & & \\\midrule
ได้ผล & 6 / 6 & 6 / 6 & 6 / 6 & 3 / 6\\
เฉลี่ย & 2.08 & 2.35 & 2.33 & 2.17\\\bottomrule
\end{tabular}
\end{table}


\chapter{ตัวอย่างข้อมูลและแนวทางตรวจซ้ำ}
\section{ตัวอย่างภาพทดสอบและสิ่งที่ระบบตอบ}
ตารางนี้แสดงภาพตัวอย่าง 4 ภาพจากชุดทดสอบ พร้อมคำบรรยายที่ระบบสร้างจริงในไฟล์ \file{results.json} โดยไม่แก้ไขข้อความ ช่อง ``คำค้น'' คือคำที่ใช้ค้นหาภาพนั้นมา ซึ่งใช้เป็นสิ่งที่คาดหวังอย่างหยาบเท่านั้น

\newcommand{\example}[6]{\noindent\begin{minipage}[t]{38mm}\vspace{0pt}\includegraphics[width=38mm,height=34mm,keepaspectratio]{../data/img/#1}\end{minipage}\hfill
\begin{minipage}[t]{\dimexpr\linewidth-42mm}\vspace{0pt}\textbf{\file{#1}}\quad คำค้น: #2\par
\textbf{ระบบบรรยาย:} #3\par
\textbf{อารมณ์ที่ระบบระบุ:} #4\par
\textbf{ตรงกับคำค้นหรือไม่:} #5\par
{\small เครดิตภาพ: #6}\end{minipage}\par\medskip}
{\centering\captionof{table}{ภาพตัวอย่าง คำบรรยายจากระบบ และเครดิตภาพ}\par}\medskip
\example{emotion_0.jpg}{คนยิ้ม (smiling face portrait)}{หญิงสาวหน้าตายิ้มแย้มแจ่มใสดูมีความสุขมาก โดยมีจักรยานสีแดงที่บรรทุกสัมภาระจอดอยู่ด้านหลังในพื้นที่กลางแจ้ง}{happy (มีความสุข) ความมั่นใจสูง}{ตรง}{``rachel smiling'' โดย .jocelyn. สัญญาอนุญาต CC BY จาก Flickr}
\example{emotion_1.jpg}{คนร้องไห้ (crying person face)}{ชายหนุ่มสวมเสื้อโปโลลายขวางนั่งอยู่บนเก้าอี้เหล็กกลางแจ้ง เขากำลังอ้าปากกว้างและทำหน้าตาเหมือนคนกำลังหาวหรือจาม}{neutral (เฉย ๆ) ความมั่นใจต่ำ}{ไม่ตรง ต้องให้คนตรวจภาพจริงว่าคำค้นหรือระบบถูก}{``My Oh Face'' โดย greggoconnell สัญญาอนุญาต CC BY จาก Flickr}
\example{emotion_2.jpg}{คนโกรธ (angry face person)}{ชายคนหนึ่งกำลังทำหน้าตาโกรธจัดใส่กล้อง คิ้วขมวดแน่นจนหน้าผากย่นและจมูกย่นเข้าหากัน ดูเหมือนเขากำลังไม่พอใจอย่างมาก}{angry (โกรธ) ความมั่นใจสูง}{ตรง}{``Angry @ The Terrorists'' โดย greggoconnell สัญญาอนุญาต CC BY จาก Flickr}
\example{env_0.jpg}{ภายในร้านกาแฟ (coffee shop interior)}{ภายในร้านมีบรรยากาศคลาสสิก มีชายคนหนึ่งนั่งนิ่งอยู่บนเก้าอี้ไม้ และอีกคนนั่งหันหลังให้บนเก้าอี้สีแดงข้างตู้ลิ้นชักไม้ขนาดใหญ่ \textbf{คำเตือน:} พื้นที่ค่อนข้างแคบ อาจเดินชนตู้ได้ง่าย}{neutral (เฉย ๆ) ความมั่นใจปานกลาง}{ตรง}{``Coffee beans for sale -- Brother Baba Budan'' โดย avlxyz สัญญาอนุญาต CC BY-SA จาก Flickr}
URL เต็มของทุกภาพอยู่ใน \file{data/credits.csv} ภาพ 4 ภาพนี้เป็นเพียงตัวอย่าง ไม่ใช่ผลประเมินทั้งชุด

\section{ข้อควรระวังในการรันทดสอบซ้ำ}
สคริปต์ \file{eval/run.py} เขียนทับ \file{results.json} และ \file{score.csv} จึงต้องสำรองไฟล์หรือรันในสำเนาโครงการก่อน การรันยังเรียกบริการภายนอกและใช้โควตา API ต้องกำหนดชุดภาพและค่าทดลองให้แน่นอนก่อนเริ่ม ส่วนการสรุปไฟล์ผลเดิมสามารถทำโดยไม่ต้องเรียกโมเดลใหม่

\section{ข้อมูลที่ยังต้องเติมก่อนรับรองการใช้งานจริง}
ยังต้องจัดทำรายงานประเมินโดยมนุษย์ ผลทดสอบกับผู้ใช้เป้าหมาย บันทึกฮาร์ดแวร์และเวอร์ชัน dependency ที่ตรึงแล้ว ผลสอบเทียบระยะจริง และนโยบายข้อมูลสำหรับบริการที่นำไปใช้งาน ไม่ควรเติมข้อมูลเหล่านี้ย้อนหลังหากไม่มีหลักฐานการดำเนินการ
\end{document}
